\chapter{HDRI}

High Dynamic Range Imaging (HDRI) is a technique used in photography and computer graphics to capture and represent a wider range of luminosity than what is possible with standard imaging techniques. HDRI allows for the representation of both very bright and very dark areas in a scene, which can be particularly useful in situations where there is a significant contrast between light and shadow. In real live, dark colors can be $10^4$ times darker than bright colors ($1:10000$), while normal images that use $8$ bits per channel can only represent a range of $1:256$. During this chapter, we will first discuss how does the human work, and then the three main stages of HDRI will be covered: recording, storage, and output.

\section{Nature}

\section{Recording}

\section{Storage}

\section{Output}

\subsection{Light Bloom}

\subsection{Real-time Tone Mapping}